% =====================================================================
% preambulo.tex: modelo fixo dos relatórios de AEDS2 (PUCRS)
%
% Para um novo trabalho: copie a pasta inteira, troque título e autor em
% main.tex e reescreva as seções em secoes/. Este arquivo não precisa mudar.
%
% Compila com pdfLaTeX (padrão do Overleaf) e com XeLaTeX. O ramo \ifPDFTeX
% escolhe as fontes de cada motor; o resto é igual.
% =====================================================================

\usepackage{iftex}
\ifPDFTeX
  \usepackage[utf8]{inputenc}
  \usepackage[T1]{fontenc}
\else
  \usepackage{fontspec}
\fi
\usepackage[brazilian]{babel}
\usepackage{amsmath}
\ifPDFTeX
  \usepackage{newtxtext,newtxmath}   % Times no texto e na matemática
\else
  % Fontes chamadas pelo nome do arquivo: funciona no XeLaTeX completo e no Tectonic.
  \setmainfont{texgyretermes}[Extension=.otf,UprightFont=*-regular,BoldFont=*-bold,
    ItalicFont=*-italic,BoldItalicFont=*-bolditalic]
  \setmonofont{texgyrecursor}[Extension=.otf,UprightFont=*-regular,BoldFont=*-bold,
    ItalicFont=*-italic,BoldItalicFont=*-bolditalic,Scale=MatchLowercase]
\fi

% ---------------------------------------------------------------- página e texto
\usepackage[a4paper,margin=2.5cm]{geometry}
\setlength{\parindent}{0pt}
\setlength{\parskip}{0.6em}
\emergencystretch=1.5em

\usepackage{enumitem}
\setlist{itemsep=0.25em,topsep=0.3em}

% ---------------------------------------------------------------- cores, figuras, gráficos
\usepackage{xcolor}
\usepackage{graphicx}
\usepackage{tikz}
\usetikzlibrary{arrows.meta}
\usepackage{pgfplots}
\pgfplotsset{compat=1.17}

% ---------------------------------------------------------------- tabelas e legendas
\usepackage{array}
\renewcommand{\arraystretch}{1.2}
\usepackage[font=small,labelfont=it,textfont=it,labelsep=endash,justification=centering]{caption}
\captionsetup{skip=6pt}

% ---------------------------------------------------------------- pseudo-código numerado
% Linhas numeradas permitem citar "linha 7" no texto e ligar cada etapa à explicação.
\usepackage{algorithm}
\usepackage{algpseudocode}
\floatname{algorithm}{Algoritmo}
\algrenewcommand\algorithmicreturn{\textbf{devolve}}
\algrenewtext{Function}[2]{\textbf{função} \textproc{#1}(#2)}
\algrenewtext{EndFunction}{\textbf{fim função}}
\algrenewtext{While}[1]{\textbf{enquanto} #1 \textbf{faça}}
\algrenewtext{EndWhile}{\textbf{fim enquanto}}
\algrenewtext{If}[1]{\textbf{se} #1 \textbf{então}}
\algrenewtext{EndIf}{\textbf{fim se}}
\algrenewtext{For}[1]{\textbf{para} #1 \textbf{faça}}
\algrenewtext{ForAll}[1]{\textbf{para cada} #1 \textbf{faça}}
\algrenewtext{EndFor}{\textbf{fim para}}

% ---------------------------------------------------------------- PDF
\usepackage[hidelinks]{hyperref}

% ---------------------------------------------------------------- macros do modelo
% \pendente{texto}: lembrete em vermelho para resolver antes de entregar.
% Antes da entrega, procure por "pendente" nos arquivos .tex e apague todos.
\newcommand{\pendente}[1]{\textcolor{red}{\textbf{[#1]}}}

% \espacofigura{altura}{indicação}: caixa tracejada que reserva o lugar de uma figura
% e diz o que desenhar. Quando a imagem existir, troque a chamada por
% \includegraphics[width=0.9\linewidth]{figuras/arquivo}.
\newcommand{\espacofigura}[2]{%
  \begin{tikzpicture}
    \node[draw=gray,dashed,rounded corners=2pt,text width=0.86\linewidth,
          minimum height=#1,align=center,inner sep=6pt,font=\small\itshape]
      {\textbf{Espaço para a figura}\\[3pt] #2};
  \end{tikzpicture}}
